\chapter{Security Considerations}\label{chap:security}

In this chapter, we record the security baseline we observed at the source commits
listed in Chapter~\ref{chap:results}. We cover network access, authentication and
authorization, payment confirmation, secrets management, IAM, the data layer, and
pod-level hardening. However, this baseline predates our later production
remediation and retest in Section~\ref{sec:testing}; where the two differ, we use
the retest as the newer operational evidence.

\section{Network-Level Access}\label{sec:sec-network}
We summarise the network paths we found open to the
cluster, based on \texttt{modules/\allowbreak network/\allowbreak locals.tf} and the boolean overrides in
\texttt{main.tf} (Section~\ref{sec:domain-context}).

\begin{table}[H]
\centering
\caption{Network access rules.}
\label{tab:security-network}
\small
\begin{tabularx}{\textwidth}{>{\raggedright\arraybackslash}p{3.4cm} X}
\toprule
\textbf{Path} & \textbf{Current rule} \\
\midrule
Kubernetes API endpoint & TCP 6443 open to \texttt{0.0.0.0/0}, not restricted to team IPs \\
Load Balancer & TCP 80 and TCP 443 both open to \texttt{0.0.0.0/0} \\
Worker node SSH & TCP 22 open to \texttt{0.0.0.0/0} (\texttt{enable\_worker\_ssh\_access = true} in \texttt{main.tf}) \\
Worker node NodePort range & TCP 30000--32767 open to \texttt{0.0.0.0/0}, bypassing the Load Balancer (\texttt{enable\_nodeport\_public\_access = true}) \\
Worker node general outbound & Unrestricted outbound TCP to \texttt{0.0.0.0/0} (\texttt{enable\_worker\_\allowbreak general\_\allowbreak outbound = true}) \\
Pod outbound & TCP 443 to \texttt{0.0.0.0/0} (\texttt{enable\_pods\_outbound\_internet = true}) \\
ArgoCD server UI & Exposed via a public NGINX ingress (\texttt{argocd/argocd-ingress.yaml}), HTTPS force-redirected \\
\bottomrule
\end{tabularx}
\end{table}

We found that the network module's optional access switches default to off;
however, the root module enables all of them. Worker nodes and pods still reach OCI services (including
OCIR image pulls) through the private Service Gateway rather than the NAT path,
matching the route tables in Appendix~\ref{sec:appendix-network}.

\section{Authentication and Authorization}\label{sec:sec-app-auth}
\begin{warningbox}[Backend Authorization at the Initial Review]
In our initial review, we found that the backend services (\texttt{auth-service}, \texttt{products-service},
\texttt{orders-service}, \texttt{payments-service}) do not check a token, header, or
role on any route. \texttt{products-service}'s \texttt{POST}, \texttt{PUT}, and
\texttt{DELETE /products} routes -- and the \texttt{ecomm-ui} BFF routes and Server
Actions that call them -- accept requests from any caller that can reach the
service.
\end{warningbox}

We found two authentication paths in the system:

\begin{itemize}[leftmargin=1.4em]
  \item \textbf{Admin authentication.} \texttt{auth-service}'s \texttt{POST
        /auth/admin/login} issues a \texttt{jose}-signed HS256 JWT (1-day expiry)
        after checking a \texttt{bcryptjs} password hash. \texttt{ecomm-ui} stores
        the token in an httpOnly \texttt{admin\_token} cookie
        (\texttt{/api/admin/login/route.ts}), and \texttt{middleware.ts} verifies
        that cookie for requests to \texttt{/admin/*} pages. The underlying API
        routes those pages call are not covered by this check.
  \item \textbf{Standard user login.} \texttt{auth-service}'s \texttt{POST
        /auth/user/login} checks the password hash and returns user data but does
        not issue a token. \texttt{ecomm-ui}'s login page notes this directly in a
        comment: \texttt{// In a real app, save token/session here}. No
        user-scoped session exists after a successful login.
\end{itemize}

\section{Payment Confirmation}\label{sec:sec-payments}
We found that \texttt{payments-service} creates a Stripe Checkout Session and a
\texttt{pending} \texttt{Order} record. When Stripe redirects the browser back to
\texttt{/checkout/success}, that page calls \texttt{POST /api/orders/confirm} with
the \texttt{orderId} and \texttt{sessionId} from the URL. The route does not call
Stripe to verify the session before updating the order; it sets \texttt{status:
"paid"} directly. The frontend's own source acknowledges this:

\begin{lstlisting}[language={}, caption={\texttt{src/app/(store)/checkout/success/page.tsx}}]
// In a real production app, you would verify the session status
// with your backend here, or use a Stripe Webhook to mark the order paid.
// For this tutorial, we will trust the redirect.
\end{lstlisting}

We also found no Stripe webhook handler in \texttt{payments-service} or
elsewhere in the repositories at the reviewed commits. Our subsequent remediation added Stripe
session validation: the production retest rejected a fabricated session ID and
closed \code{BUG-018} (Section~\ref{sec:testing}).

\section{Secrets Management}\label{sec:sec-secrets}
We found two different secrets-management patterns across the repository:

\begin{itemize}[leftmargin=1.4em]
  \item \textbf{Automated deployment (\texttt{deploy.yml}).} \texttt{JWT\_SECRET},
        \texttt{STRIPE\_SECRET\_KEY}, \texttt{ADMIN\_EMAIL}, and
        \texttt{ADMIN\_PASSWORD} are read from encrypted GitHub Actions secrets and
        written into a Kubernetes Secret (\texttt{app-secrets}) at deploy time.
        The same \texttt{ADMIN\_EMAIL} secret is also passed to cert-manager's
        \texttt{letsencrypt-prod} \texttt{ClusterIssuer} as the ACME registration
        email (Section~\ref{sec:bootstrap-pipeline}).
        \texttt{DEPLOYMENT.md} documents an equivalent manual path using a
        gitignored \texttt{.env} file. Neither commits real secret values to Git.
  \item \textbf{Plaintext values in \texttt{values.yaml}.}
        \texttt{helm/ecommerce-app/values.yaml} also contains a plaintext
        \texttt{secrets:} block with a JWT secret, a Stripe secret key in live key
        format, and a default admin email/password. No Helm template references
        \texttt{.Values.secrets}, so the block does not affect what gets deployed.
        It should still be removed and the associated credentials rotated, since it
        remains in the repository's Git history.
\end{itemize}

Moreover, we found that OCI Vault integration is referenced in \texttt{DEPLOYMENT.md}, which lists three
GitHub secret names implying Vault-stored OCIDs
(\texttt{VAULT\_OCIR\_AUTH\_TOKEN\_OCID}, \texttt{VAULT\_STRIPE\_SECRET\_KEY\_OCID},
\texttt{VAULT\_JWT\_SECRET\_OCID}). None of these names appear in any Terraform file
or GitHub Actions workflow, and no Vault resource exists in
\texttt{Terraform-code}; secrets are managed through Kubernetes Secrets instead.

\section{IAM}\label{sec:sec-iam}
In our IAM review, we found that \texttt{main.tf} contains a commented-out dynamic group and IAM policy for OKE
worker access to images and secrets, with an inline note that the team did not have
tenancy-level access to create dynamic groups. The current deployment uses a
static, permanent OCI Auth Token (\texttt{ocir-secret}, a Kubernetes
\texttt{docker-registry} Secret) instead. An earlier version generated this token
from the OCI CLI's short-lived \texttt{access-token get} command, which expired and
took down image pulls cluster-wide with \texttt{ImagePullBackOff} once it did; the
permanent token avoids that failure mode.

\section{Data Layer}\label{sec:sec-data}
We observed that MongoDB is deployed as a single replica
(\texttt{helm/ecommerce-app/\allowbreak templates/\allowbreak mongodb.yaml}) with no
authentication configured. Any pod that can reach the \texttt{mongodb} ClusterIP
Service in the namespace can read or write every collection without credentials, and
the single-replica deployment is a shared dependency for all five workloads with no
built-in redundancy.

\section{Pod-Level Hardening}\label{sec:sec-pod}
However, we also confirmed a consistent pod-hardening measure: every workload's
Deployment sets \texttt{securityContext: runAsNonRoot: true} with an
explicit non-root UID (1000 for the four backend services, 1001 for
\texttt{ecomm-ui}), matching each Dockerfile's own \texttt{USER} directive, applied
consistently across all five workloads.
